\section{Related Work}
\label{sec:related}

\subsection{Ocular-Artifact Correction}
Classical approaches include EOG regression, blind source separation, and
artifact subspace reconstruction. Regression explicitly maps EOG into EEG and
is attractive when EOG remains available at deployment, although transfer
coefficients may change across participants and sessions. ICA separates latent
components but requires sufficiently stable decomposition and a component
selection rule. SGEYESUB provides a careful participant-level evaluation of
ocular correction and illustrates the need to measure both attenuation and
neural preservation \cite{kobler2020sgeyesub}.

Learned denoisers instead infer a correction from population data. Paired
benchmarks such as EEGdenoiseNet enabled controlled comparisons of convolutional,
recurrent, and generative estimators \cite{zhang2021eegdenoisenet}. These
methods establish method-level denoising efficacy, but do not by themselves
test whether a disjoint calibration from an unseen participant adds value.

\subsection{Diffusion Denoising}
Diffusion models learn a reverse process from progressively corrupted states
\cite{ho2020ddpm}. EEGDfus and D4PM apply diffusion ideas to EEG artifact
removal under paired or mixture assumptions \cite{eegdfus2024,d4pm2025}.
Their objectives, conditioning information, and benchmark protocols differ
from ours. We do not treat those differences as shortcomings: they answer
whether a denoising method works on their target benchmarks, whereas our
experiment isolates matching, population, and mismatched-subject contexts
inside one frozen EOG-guided diffusion pipeline.

\subsection{Calibration and Subject Adaptation}
EEG variability motivates alignment, normalization, domain adaptation, and
participant-conditioned representations. A known participant identifier or a
closed-set embedding is not directly available for an unseen user. Our setting
instead derives the context from early synchronized EEG+EOG and applies it to
later data. The claim is correspondingly narrow: matching support improves this
fixed estimator relative to two operator controls. We do not claim that
calibration is unique to diffusion, and we retain linear and deterministic
comparators to expose that boundary.
